% ECONOMICS_PROBLEM: {"id": "F13-636", "title": "Services and digital trade barriers", "jel_code": "F13", "source_id": "OP-0089"}
% !TeX program = xelatex
\documentclass[11pt,a4paper]{article}
\usepackage[margin=23mm]{geometry}
\usepackage{fontspec}
\setmainfont{Latin Modern Roman}
\usepackage{amsmath,amssymb,mathtools}
\usepackage{xcolor,enumitem,booktabs,longtable,array,xurl,hyperref}
\definecolor{muted}{HTML}{53636C}
\hypersetup{colorlinks=true,urlcolor=blue,linkcolor=blue}
\setlength{\parindent}{0pt}
\setlength{\parskip}{5pt}
\newcommand{\R}{\mathbb R}
\newcommand{\E}{\mathbb E}
\newcommand{\N}{\mathbb N}
\newcommand{\ind}{\mathbf 1}
\newcommand{\Var}{\operatorname{Var}}
\newcommand{\Cov}{\operatorname{Cov}}
\newcommand{\supp}{\operatorname{supp}}
\DeclareMathOperator*{\argmin}{arg\,min}
\DeclareMathOperator*{\argmax}{arg\,max}
\newcommand{\entrymeta}[3]{{\sffamily\small\color{muted}\textbf{\texttt{#1}}\quad|\quad #2\par #3\par}}
\newcommand{\Problemref}[1]{\texttt{#1}}
\newcommand{\JELref}[2]{\texttt{#2} (JEL \texttt{#1})}
\begin{document}
\section*{Services and digital trade barriers}
\textbf{Problem ID:} \texttt{F13-636}\quad \textbf{JEL:} \texttt{F13}\par
% BEGIN ECONOMICS STATEMENT
\label{problem:OP-0089}
\label{atlas:prob:T9}

\noindent\textbf{Atlas ID:} T9; \textbf{original number:} 89; \textbf{field:} International trade. \textbf{Classification:} Editorial research formulation (theoretical/empirical); not a certified unsolved theorem.

\subsection*{Definitions and assumptions}
Fix a services sector, origin and destination markets, household types and a structural model $\theta$. A regulatory regime $r$ restricts licensing, establishment, cross-border data transfers or remote delivery through explicitly defined feasible firm choices. Firms set prices, invest and supply services; consumers choose services within their budgets; a stated equilibrium selection clears markets. A comparison regime with an ad valorem tariff $t\ge0$ applies a factor $1+t$ to a defined cross-border transaction value and specifies revenue recycling. Let $Q_\theta(r)$ be service quantity, $W_{h,\theta}(r)$ household welfare and $N_\theta(r)$ firm entry.

\subsection*{Formal research question}
For a baseline $r^0$ and reform $r^1$, ask whether a tariff $t^*$ exists and is unique such that
\[
 Q_\theta(r^1)=Q_\theta(r^0,t^*),
\]
where $(r^0,t)$ denotes the baseline regulation with the specified tariff. Compare this quantity-equivalent value with tariffs matching household welfare or entry. Identify the set of equivalent tariffs from transaction, price and regulatory data when $\theta$ is partially identified. Characterize conditions under which one scalar tariff reproduces a vector of quantity, entry and welfare effects, and construct counterexamples when regulations act through fixed costs, quality, privacy or market access rather than marginal transaction costs.

\subsection*{Interpretation and scope}
The comparison baseline and direction of the tariff adjustment must be specified. A liberalizing reform may require a subsidy, so restricting $t\ge0$ can prevent existence; an alternative domain $t>-1$ should be stated if permitted. Continuity and strict monotonicity of the chosen outcome in $t$ support uniqueness, while a matching root also requires the outcome lie in its attainable range. Digitally delivered services do not imply zero production, compliance or matching costs. Platform data may omit informal transactions or misclassify purchaser location. A numerical barrier equivalent is model- and outcome-dependent; it cannot automatically be interpreted as a welfare-equivalent regulatory burden.

\subsection*{Source audit and status}
[S1] concerns the Internet and historical services trade; [S2] concerns distance and offshoring. They motivate the subject but do not identify contemporary licensing or data-policy tariff equivalents. The atlas’s ``digital-trade literature'' lacks author, title and date and remains unresolved. The question is an editorial measurement and identification agenda, without certification of a currently open mathematical conjecture.

\subsection*{References}

\begin{enumerate}[label={[S\arabic*]},leftmargin=*]

\item Caroline Freund and Diana Weinhold (2002). \emph{The Internet and International Trade in Services}. American Economic Review 92(2), 236–240. \url{https://doi.org/10.1257/000282802320189320}.
\textit{Locator and role:} Publisher or author abstract / bibliographic record; Historical association between Internet use and services trade. Provides background or a benchmark result; does not establish that the editorial question remains unresolved in 2026.

\item Keith Head, Thierry Mayer, and John Ries (2009). \emph{How Remote Is the Offshoring Threat?}. European Economic Review 53(4), 429–444. \url{https://www.sciencedirect.com/science/article/pii/S0014292108000822}.
\textit{Locator and role:} Publisher or author abstract / bibliographic record; Distance and offshoring in disaggregated services trade. Provides background or a benchmark result; does not establish that the editorial question remains unresolved in 2026.

\item \textbf{Unresolved original citation:} ``digital-trade literature.'' No author, title or year was supplied, so no unique source can be assigned.

\end{enumerate}

\subsection*{Common conventions: Source mathematical conventions}

Unless an entry states otherwise, agent and firm sets are finite. State and action spaces are measurable, strategies and outcome maps are measurable, probability laws are specified, and expectations exist for the targets under discussion. Infinite-horizon discount factors lie in $(0,1)$ and discounted objectives are finite. Norms, tolerances and complexity input sizes must be specified for computational comparisons. Constants or primitives that appear locally are inputs, not empirical facts asserted by this catalogue.

\textbf{Identification.} A statistical model $m\in\mathcal M$ implies an observable law $P_m$ and target $\theta(m)$. At law $P$, its identified set is
\[
\Theta(P;\mathcal M)=\{\theta(m):m\in\mathcal M,\ P_m=P\}.
\]
Point identification holds when this set is a singleton, or equivalently when $P_{m_1}=P_{m_2}$ implies $\theta(m_1)=\theta(m_2)$ for all compatible models. Sharp bounds mean bounds attained, or approached when the boundary is not attained, by compatible models. Writing this set is a definition, not a solution: the substantive task is to establish informative restrictions and derive or estimate the set. An unrestricted-confounding model may give only trivial bounds.

\textbf{Inference.} Identification, estimability and valid uncertainty quantification are separate questions. Fix $\alpha\in(0,1)$. A confidence procedure $C_n$ over a declared law class $\mathcal P$ has asymptotic uniform coverage at level $1-\alpha$ when
\[
\liminf_{n\to\infty}\inf_{P\in\mathcal P}\Pr_P\{\theta(P)\in C_n\}\geq1-\alpha.
\]
For set-identified targets the coverage event must be defined separately, for example coverage of the full identified set. If an entry uses identification-indexed models instead of uniquely indexed laws, the same coverage convention is applied over those models. Pointwise limits do not establish uniform validity, and asymptotic validity does not guarantee finite-sample precision. Sample sizes, dependence structures and weak-identification sequences are part of each application.

\textbf{Counterfactuals.} $Y_i(a)$ is a potential outcome under a specified intervention, including its timing, implementation and versions. It is not $Y_i$ conditional on voluntarily choosing $a$. A full intervention vector permits interference; shorthand scalar treatment notation does not silently impose no interference. Assignment, selection, attrition, anticipation and measurement must be specified. A claim about an instrument's local effect is not automatically a population policy effect.

\textbf{Economic equilibrium.} A counterfactual model includes best responses, constraints, feasibility, market clearing, state transitions and expectations consistent with the stated information. When several equilibria exist, either a declared selector is an input or the target is set-valued. Comparisons across policy regimes must use the stated selection convention. Reduced-form relationships are not presumed invariant after an equilibrium-changing intervention.

\textbf{Optimization and welfare.} Optimization is over the stated feasible set. If attainment is not established, $\sup$ or $\inf$ and $\varepsilon$-optimal choices replace an asserted maximizer or minimizer. For example, with value $V=\sup_{a\in\mathcal A}W(a)<\infty$, an $\varepsilon$-optimal set is $\{a\in\mathcal A:W(a)\geq V-\varepsilon\}$ for $\varepsilon>0$. Benefits, costs, risk and health or environmental outcomes can be aggregated only using declared units and conversion weights. Interpersonal utility weights, the treatment of fiscal transfers and foreign profits, discounting, and the behavioral welfare criterion are explicit inputs. A comparative-static derivative additionally requires existence and appropriate differentiability of the selected counterfactual.

\textbf{Sources.} Local labels [S1], [S2], and so forth restart in every question. Locators identify the relevant abstract, model, section or result at the level actually checked; a bibliographic or abstract-level verification is not represented as inspection of an entire proof. Working papers are distinguished from later publications and changing versions. Source summaries are paraphrased. All URL checks and editorial formulations refer to the audit date stated above.
% END ECONOMICS STATEMENT
\end{document}
